% ============================================================
%  Chapter 4 — Polynomials
% ============================================================
\section{Polynomials}
\label{sec:polynomials}

Polynomials over a field $\F$ are the central objects of study when we analyse
eigenvalues, minimal polynomials, and the structure theorem for linear
operators.  This chapter collects the essential algebraic facts.

% ------------------------------------------------------------
\begin{defbox}[Definition --- Polynomial over $\F$]
A \textbf{polynomial} over $\F$ is an expression
\[
    p(x) = a_n x^n + a_{n-1}x^{n-1} + \cdots + a_1 x + a_0,
    \quad a_i\in\F,\; n\in\N_0.
\]
The set of all polynomials over $\F$ is denoted $\F[x]$.
\end{defbox}

% ------------------------------------------------------------
\begin{defbox}[Definition --- Degree, Leading Term, and Related Notions]
For a non-zero polynomial $p(x) = a_n x^n + \cdots + a_0 \in \F[x]$:
\begin{itemize}
    \item $\deg(p)$: the \emph{degree}, i.e.\ the largest $n$ s.t.\ $a_n\neq 0$.
    \item \emph{Leading term}: $a_n x^n$ where $n=\deg(p)$.
    \item \emph{Leading coefficient}: $a_n$.
    \item \emph{Constant coefficient}: $a_0$.
    \item \emph{Monic}: leading coefficient equals $1$.
    \item \emph{Constant polynomial}: $\deg(p)=0$, i.e.\ $p(x)=a_0\neq 0$.
\end{itemize}
By convention: $\deg(0) = -\infty$ and $x^0 = 1$.
\end{defbox}

% ------------------------------------------------------------
\begin{defbox}[Definition --- Divisibility of Polynomials]
We say $d(x)$ \textbf{divides} $p(x)$, written $d(x)\mid p(x)$, if there
exists $q(x)\in\F[x]$ s.t.\ $p(x)=d(x)q(x)$.
\end{defbox}

% ------------------------------------------------------------
\begin{propbox}[Proposition --- Basic Divisibility Properties]
For $f,g,h\in\F[x]$ and $c\in\F\setminus\{0\}$:
\begin{enumerate}
    \item $f(x)\mid 0$; \quad $c\mid f(x)$; \quad $cf(x)\mid f(x)$.
    \item If $f\mid g$ and $g\mid h$, then $f\mid h$.
    \item If $f\mid g$, then $f\mid g\cdot h$.
    \item If $f\mid g$ and $f\mid h$, then $f\mid (g+h)$.
    \item If $f(x)\mid c$, then $\deg(f)=0$.
    \item If $f\mid g$ and $g\mid f$, then $\exists\, d\in\F\setminus\{0\}$ s.t.\ $f=dg$.
    \item $f\mid g \iff cf\mid g \iff f\mid cg$.
\end{enumerate}
\end{propbox}

% ------------------------------------------------------------
\begin{thmbox}[Theorem --- Greatest Common Divisor of Polynomials]
Let $f,g\in\F[x]$, not both zero.  Then there exists a \textbf{unique monic}
polynomial $h(x)$ of highest degree that divides both $f$ and $g$.  Moreover:
\begin{itemize}
    \item Every common divisor of $f$ and $g$ divides $h$.
    \item There exist $r,t\in\F[x]$ s.t.\ $h = rf + tg$ (Bézout's identity).
\end{itemize}
We write $\gcd(f,g)=h$, and call $f,g$ \textbf{coprime} if $\gcd(f,g)=1$.
\end{thmbox}

% ------------------------------------------------------------
\begin{thmbox}[Theorem --- Long Division of Polynomials]
For $p,q\in\F[x]$ with $q\neq 0$, there exist unique $c,r\in\F[x]$ s.t.
\[
    p(x) = c(x)q(x) + r(x), \quad \deg(r) < \deg(q) \text{ or } r=0.
\]
\end{thmbox}

% ------------------------------------------------------------
\begin{thmbox}[Theorem --- Remainder Theorem]
For any $p(x)\in\F[x]$ and $\alpha\in\F$,
\[
    p(x) = (x-\alpha)\,c(x) + p(\alpha).
\]
\end{thmbox}

% ------------------------------------------------------------
\begin{corbox}[Corollary --- Factor Theorem]
$(x-\alpha)\mid p(x)$ if and only if $p(\alpha)=0$.
\end{corbox}

% ------------------------------------------------------------
\begin{thmbox}[Theorem --- Coprime Divisibility (Lemma)]
If $f\mid gh$ and $\gcd(f,g)=1$, then $f\mid h$.
\end{thmbox}

% ------------------------------------------------------------
\begin{corbox}[Corollary --- Coprime Factor Product]
If $f\mid h$ and $g\mid h$ with $\gcd(f,g)=1$, then $fg\mid h$.
\end{corbox}

% ------------------------------------------------------------
\begin{thmbox}[Theorem --- Unique Factorisation of Polynomials]
Let $f(x)\in\F[x]$, $f\neq 0$, $\deg(f)\geq 1$.  Then there exist irreducible
monic polynomials $p_1,\ldots,p_r\in\F[x]$ and a scalar $c\in\F\setminus\{0\}$
s.t.
\[
    f(x) = c\,p_1(x)^{n_1}\cdots p_r(x)^{n_r}.
\]
The polynomials $p_1,\ldots,p_r$ are unique up to permutation; the exponents
$n_1,\ldots,n_r$ and the scalar $c$ are uniquely determined.
\end{thmbox}

% ------------------------------------------------------------
\begin{thmbox}[Theorem --- Fundamental Theorem of Algebra]
$\C$ is \emph{algebraically closed}: every non-constant polynomial in $\C[x]$
has at least one root in $\C$.

Equivalently, every $f(x)\in\C[x]$ of degree $n\geq 1$ factors completely over
$\C$ into $n$ linear factors (counting multiplicity).
\end{thmbox}

% ------------------------------------------------------------
\begin{thmbox}[Theorem --- Gauss' Rational Root Criterion]
Let $f(x) = a_n x^n + \cdots + a_0\in\Z[x]$.  If $r/s$ (with $\gcd(r,s)=1$)
is a rational root of $f$, then $s\mid a_n$ and $r\mid a_0$.
\end{thmbox}

% ------------------------------------------------------------
\begin{thmbox}[Theorem --- Eisenstein's Irreducibility Criterion]
Let $f(x) = x^n + a_{n-1}x^{n-1} + \cdots + a_0\in\Z[x]$ be monic.  If there
exists a prime $p$ s.t.
\begin{enumerate}
    \item $p\mid a_i$ for all $i\in\{0,1,\ldots,n-1\}$, and
    \item $p^2\nmid a_0$,
\end{enumerate}
then $f$ is irreducible over $\Q[x]$.
\end{thmbox}

% ------------------------------------------------------------
\begin{thmbox}[Theorem --- Vieta's Formulae]
Let $p(x) = a_n x^n + \cdots + a_0 = a_n(x-x_1)\cdots(x-x_n)\in\F[x]$.
Then the roots $x_1,\ldots,x_n$ satisfy:
\begin{align*}
    \sum_{i} x_i
        &= -\,\frac{a_{n-1}}{a_n}, \\
    \sum_{i<j} x_i x_j
        &= \phantom{-}\frac{a_{n-2}}{a_n}, \\
    \sum_{i<j<k} x_i x_j x_k
        &= -\,\frac{a_{n-3}}{a_n}, \\
        &\;\;\vdots \\
    \prod_{i} x_i
        &= (-1)^n\frac{a_0}{a_n}.
\end{align*}
\end{thmbox}
